\section{Related Work}
\label{sec:related_work}

\noindent\textbf{Vision-Language-Action Models.} Recent work has developed generalist robot policies trained on large and diverse robotic datasets \cite{embodimentcollaboration2025openxembodimentroboticlearning, khazatsky2025droidlargescaleinthewildrobot, ebert2021bridgedataboostinggeneralization, pmlr-v229-walke23a, fang2023rh20tcomprehensiveroboticdataset, dasari2020robonetlargescalemultirobotlearning, wu2025robomindbenchmarkmultiembodimentintelligence, pmlr-v164-jang22a, galaxea2025, agibotworldcontributors2025agibotworldcolosseolargescale}. Among these, Vision--Language--Action models (VLAs) have emerged as a particularly promising paradigm \cite{brohan2023rt1roboticstransformerrealworld, embodimentcollaboration2025openxembodimentroboticlearning, kim2024openvlaopensourcevisionlanguageactionmodel, octomodelteam2024octoopensourcegeneralistrobot, black2024pi0visionlanguageactionflowmodel, li2024cogactfoundationalvisionlanguageactionmodel, geminiroboticsteam2025geminiroboticsbringingai, liu2025rdt1bdiffusionfoundationmodel, brohan2023rt2visionlanguageactionmodelstransfer, nvidia2025gr00tn1openfoundation, hancock2025actionslanguagefinetuningvlms, pertsch2025fastefficientactiontokenization, wang2024scalingproprioceptivevisuallearningheterogeneous, zheng2025tracevlavisualtraceprompting, Zhao_2025_CVPR}. These models are typically fine-tuned from Vision--Language models (VLMs) pre-trained on internet-scale data. The pre-trained representations equip VLAs with the ability to follow natural language instructions and understand diverse visual cues. As a result, VLAs often not only perform well within their training distribution, but also show promise for generalization to out-of-distribution environments \cite{intelligence2025pi05visionlanguageactionmodelopenworld, black2024pi0visionlanguageactionflowmodel, kim2024openvlaopensourcevisionlanguageactionmodel, gao2024efficientdatacollectionrobotic, zha2025guidingdatacollectionfactored, hu2025datascalinglawsimitation}. Nevertheless, despite being trained on data spanning multiple robot embodiments, existing VLAs do not generalize zero-shot to novel embodiments without extensive fine-tuning (Section~\ref{subsec:zero_shot}). In this work, we introduce a new action representation that significantly improves VLA's embodiment generalization capability.

\vspace{5pt}
\noindent\textbf{Cross-Embodiment Policy Learning.} Large robot policies are typically trained on heterogeneous, cross-embodiment datasets, and a common strategy for handling varying state and action spaces is to manually define a unified representation. The most straightforward approach pads states and actions to the maximum dimensionality or uses shared action spaces like end-effector pose \cite{black2024pi0visionlanguageactionflowmodel, kim2024openvlaopensourcevisionlanguageactionmodel, zheng2025xvlasoftpromptedtransformerscalable}. Other works learn embodiment-specific state or action projectors, or employ separate action heads for different robot types \cite{wang2024scalingproprioceptivevisuallearningheterogeneous, doshi2024scalingcrossembodiedlearningpolicy, octomodelteam2024octoopensourcegeneralistrobot, yang2024pushinglimitscrossembodimentlearning, nvidia2025gr00tn1openfoundation}. A parallel line of work learns shared latent actions or skills across embodiments \cite{ye2025latentactionpretrainingvideos, bu2025univlalearningacttaskcentric, Zheng_2025_CVPR, Chen_2025_ICCV, xu2023xskillcrossembodimentskill, wang2024crossembodimentrobotmanipulationskill}. More recently, human hand motion has been used as a unified action space bridging human and robot data \cite{luo2026beingh05scalinghumancentricrobot,yang2025egovlalearningvisionlanguageactionmodels}, enabling task generalization to seen embodiments. While effective, these methods typically require embodiment-specific fine-tuning to adapt to new embodiments, whereas our method achieves \emph{zero-shot embodiment transfer} without any additional training.

Another line of work facilitates cross-embodiment learning by conditioning robot policies on embodiment features. Early approaches infer such features from kinematic structure \cite{huang2020policycontrolallshared, kurin2021bodycagerolemorphology, gupta2022metamorphlearninguniversalcontrollers, patel2024getzerographembodimenttransformer, bohlinger2025policyrunallendtoend}, and several works focus on human-to-robot transfer by exploiting structural similarity \cite{chen2024roviaugrobotviewpointaugmentation, lum2025crossinghumanrobotembodimentgap, lepert2025phantomtrainingrobotsrobots, ren2025motiontracksunifiedrepresentation}. For example, X-VLA \cite{zheng2025xvlasoftpromptedtransformerscalable} conditions policies on per-embodiment soft prompts, and DexVLA \cite{wen2025dexvlavisionlanguagemodelplugin} introduces an explicit embodiment-specific training stage. RDT2 \cite{liu2026rdt2exploringscalinglimit} similarly reports zero-shot embodiment transfer, but does so under the assumption that new robots share the same gripper design and wrist-mounted camera configuration as those used during data collection, effectively constraining embodiment variation to matched hardware interfaces rather than arbitrary morphologies. In parallel, a broader class of approaches infers embodiment implicitly from long interaction histories, as commonly explored in navigation and embodied control \cite{hirose2025omnivlaomnimodalvisionlanguageactionmodel, shah2023vintfoundationmodelvisual, yang2024pushinglimitscrossembodimentlearning, shah2023gnmgeneralnavigationmodel}. Taken together, these methods primarily model embodiment variation within the training distribution—through prompts, explicit embodiment conditioning, hardware alignment, or historical context—and typically require additional assumptions or adaptation when deployed on previously unseen embodiments.

By contrast, LAP is motivated by the complementary goal of true zero-shot embodiment transfer. Rather than introducing explicit embodiment-conditioning mechanisms or relying on matched hardware interfaces across robots (e.g., identical grippers and wrist-mounted cameras), LAP-3B demonstrates that cross-embodiment generalization can emerge directly from large-scale VLA pre-training with appropriately structured action representations. Without any embodiment-specific design or embodiment-aligned data collection assumptions, LAP-3B, to the best of our knowledge, is the first VLA to achieve substantial \emph{zero-shot generalization to previously unseen embodiments}, while also enabling efficient adaptation to more challenging tasks when additional data are available.



\vspace{5pt}
\noindent\textbf{Action Representations for VLAs.}  How actions are represented determines how effectively a pre-trained VLM can be adapted for motor control  \cite{lee2025molmoactactionreasoningmodels, pertsch2025fastefficientactiontokenization, liu2025fasterefficientautoregressivevision, goyal2025vla0buildingstateoftheartvlas, grover2025enhancinggeneralizationvisionlanguageactionmodels, niu2024llarvavisionactioninstructiontuning, ye2025latentactionpretrainingvideos, hancock2025actionslanguagefinetuningvlms}. Early approaches represent actions as code or language-level sub-tasks \cite{ahn2022icanisay, driess2023palmeembodiedmultimodallanguage, liang2023codepolicieslanguagemodel, zha2024distillingretrievinggeneralizableknowledge, belkhale2024rthactionhierarchiesusing}. These formulations are closely related in spirit to our use of linguistic supervision, but they typically operate at coarse temporal abstraction and rely on separate low-level controllers to execute each sub-task. In contrast, LAP uses \emph{fine-grained} language-actions parsed directly from continuous end-effector motion and trains the backbone to model these motor-level effects, yielding representations tailored for low-level control. More recent works instead encode actions as 2D keypoints or visual traces \cite{liu2024mokaopenworldroboticmanipulation, niu2024llarvavisionactioninstructiontuning, lee2025molmoactactionreasoningmodels, zheng2025tracevlavisualtraceprompting, dipalo2024keypointactiontokensenable, huang2024rekepspatiotemporalreasoningrelational}, which enhance grounding but are typically still used as intermediate representations rather than direct control.

\begin{wraptable}{r}{0.6\columnwidth}
\vspace{-6pt}
\centering
\scriptsize
\setlength{\tabcolsep}{3pt}
\renewcommand{\arraystretch}{0.9}
\begin{tabular}{lccc|c}
\hline
\vspace{2pt}
 & 
\citet{grover2025enhancinggeneralizationvisionlanguageactionmodels} & 
VLA-0 \cite{goyal2025vla0buildingstateoftheartvlas} & 
VLM2VLA \cite{hancock2025actionslanguagefinetuningvlms} & 
\textbf{LAP} \\
\hline
Large-Scale Training  & \cmark & \xmark & \xmark & \cmark \\
Inference Frequency  & $\sim$0.5Hz & 4Hz & 0.2Hz & \textbf{25Hz} \\
Cross-Embodiment     & \xmark & \xmark & \xmark & \cmark \\
\hline
\end{tabular}
\caption{Comparison of VLAs that use (variants of) language-actions along key design and performance axes. LAP uniquely enables cross-embodiment generalization through large-scale pre-training and language-action representation.}
\vspace{-8pt}
\label{tab:method_comparison_axes_transposed}
\end{wraptable}


Another line of work maps actions to unused or non-linguistic tokens in VLM vocabularies \cite{brohan2023rt2visionlanguageactionmodelstransfer, kim2024openvlaopensourcevisionlanguageactionmodel, pertsch2025fastefficientactiontokenization, liu2025fasterefficientautoregressivevision}, which can induce distributional misalignment and degrade pre-trained VLM representations \cite{hancock2025actionslanguagefinetuningvlms}. More recently, approaches have begun representing actions directly in natural language \cite{hancock2025actionslanguagefinetuningvlms} or as raw string tokens \cite{goyal2025vla0buildingstateoftheartvlas, grover2025enhancinggeneralizationvisionlanguageactionmodels}, yielding improved semantic grounding, language following, and task-level generalization. As summarized in Table~\ref{tab:method_comparison_axes_transposed}, LAP builds on this emerging paradigm with two key distinctions: it introduces LAP-3B as a \emph{practical} language-action VLA for real-time control, and demonstrates that language-action representations can scale to large robot pre-training corpora—thereby enabling zero-shot cross-embodiment transfer.

